\chapter{Halbgruppenbeweis der Potenzhalbgruppe}
\hypertarget{power.proofs}{}
\noindent\PowerMainLink{}\enspace\textperiodcentered\enspace\PowerReadingLink{}
\par\medskip
\noindent\textbf{Definition der Potenzhalbgruppe.}
Sei \((A,\star)\) eine Halbgruppe. Die Menge ihrer nichtleeren Teilmengen ist
\[
 \PowNE{A}=\{X\in\mathcal P(A)\mid X\neq\varnothing\}.
\]
Die Multiplikation ist in
\FormulaRefAuto{PowerSemigroupProductDef}[def] definiert durch
\[
 \PowerProductDefinitionFormula
\]
Die Halbgruppe \((\PowNE{A},\star_{\mathcal P})\) heißt
\emph{Potenzhalbgruppe} von \((A,\star)\).

Die Mehrfachindex- und Operatorschreibweise ist in
\FormulaRefAuto{B27TermImageDef}[def] aus Band~3 eingeführt.
Die Existenz der Termbildmengen sichert
\FormulaRefAuto{TermImageSetUniqueExists}[thm].
Die Anwendung auf das Mengenprodukt ist bei
\FormulaRefAuto{PowerSemigroupProductDef}[def] in Band~28 erläutert.

\clearpage

\PowerInnerProofStatement{PowerSemigroupProductClosure}
\begin{tabproofwide}
  \proofstepwidestar[1]{\SemigroupAlg{A}{\star}}{\rA}
  \proofstepwidestar[2]{X,Y\in\PowNE{A}}{\rA}
  \proofstepwidestar[1,2]{%
    \hypertarget{power.proof.image-criteria}{}%
    \hypertarget{power.proof.PowerSemigroupProductTermImage}{}%
    X\star_{\mathcal P}Y=\{x\star y\mid x\in X,y\in Y\}}{%
    \FormulaRefAuto{PowerSemigroupProductDef}[def]{1,2,2}}
  \proofstepwidestar[1,2]{%
    \hypertarget{power.proof.PowerSemigroupProductNonempty}{}%
    X\star_{\mathcal P}Y\in\PowNE{A}}{%
    \rIEBack{3,\FormulaRefAuto{BinaryOperationImageNonemptyPowerset}[thm]{%
    \FormulaRefAuto{SemigroupBinaryOperationAxiom}[ax]{1},2}}}
\end{tabproofwide}

\PowerInnerProofStatement{PowerSemigroupCarrierElement}
\begin{tabproofwide}
  \proofstepwidestar[1]{X\in\PowNE{A}}{\rA}
  \proofstepwidestar[2]{x\in X}{\rA}
  \proofstepwidestar[1,2]{x\in A}{%
    \FormulaRefAuto{A\subseteq B,\,x\in A\vdash x\in B}[thm]{
    \rAEa{\FormulaRefAuto{NonemptyPowersetMembership}[thm]{1}},2}}
\end{tabproofwide}

\PowerInnerProofStatement{PowerSemigroupProductContains}
\begin{tabproofwide}
  \proofstepwidestar[1]{\SemigroupAlg{A}{\star}}{\rA}
  \proofstepwidestar[2]{X,Y\in\PowNE{A}}{\rA}
  \proofstepwidestar[3]{x\in X\land y\in Y}{\rA}
  \proofstepwidestar[3]{%
    \exists u\in X\,\exists v\in Y\;(x\star y=u\star v)}{%
    \rEIStar{3,3,\rII}}
  \proofstepwidestar[1,2,3]{x\star y\in X\star_{\mathcal P}Y}{%
    \rIEBack{\FormulaRefAuto{PowerSemigroupProductDef}[def]{1,2,2},%
    \FormulaRefAuto{B27TermImageMembership}[thm]{4}}}
\end{tabproofwide}

\PowerInnerProofStatement{PowerSemigroupProductSubset}
\begin{tabproofwide}
  \proofstepwidestar[1]{\SemigroupAlg{A}{\star}}{\rA}
  \proofstepwidestar[2]{X,Y\in\PowNE{A}}{\rA}
  \proofstepwidestar[1,2]{X\star_{\mathcal P}Y\subseteq A}{%
    \rAEa{\FormulaRefAuto{NonemptyPowersetMembership}[thm]{%
    \FormulaRefAuto{PowerSemigroupProductClosure}[thm]{1,2,2}}}}
\end{tabproofwide}

\par\Needspace{23\baselineskip}
\PowerInnerProofStatement{PowerSemigroupLeftOuterProductRepresentation}
\begin{tabproofwide}
  \proofstepwidestar[1]{\SemigroupAlg{A}{\star}}{\rA}
  \proofstepwidestar[2]{X,Y\in\PowNE{A}}{\rA}
  \proofstepwidestar[3]{Z\in\PowNE{A}}{\rA}
  \proofstepwidestar[1,2]{X\star_{\mathcal P}Y\in\PowNE{A}}{%
    \FormulaRefAuto{PowerSemigroupProductClosure}[thm]{1,2,2}}
  \proofstepwidestar[1,2,3]{\begin{aligned}[t]
    &(X\star_{\mathcal P}Y)\star_{\mathcal P}Z\\[-2pt]
    &\quad=\{u\star z\mid u\in X\star_{\mathcal P}Y,z\in Z\}
  \end{aligned}}{%
    \FormulaRefAuto{PowerSemigroupProductDef}[def]{1,4,3}}
\end{tabproofwide}

\par\Needspace{23\baselineskip}
\PowerInnerProofStatement{PowerSemigroupRightOuterProductRepresentation}
\begin{tabproofwide}
  \proofstepwidestar[1]{\SemigroupAlg{A}{\star}}{\rA}
  \proofstepwidestar[2]{X\in\PowNE{A}}{\rA}
  \proofstepwidestar[3]{Y,Z\in\PowNE{A}}{\rA}
  \proofstepwidestar[1,3]{Y\star_{\mathcal P}Z\in\PowNE{A}}{%
    \FormulaRefAuto{PowerSemigroupProductClosure}[thm]{1,3,3}}
  \proofstepwidestar[1,2,3]{\begin{aligned}[t]
    &\{x\star v\mid x\in X,v\in Y\star_{\mathcal P}Z\}\\[-2pt]
    &\quad=X\star_{\mathcal P}(Y\star_{\mathcal P}Z)
  \end{aligned}}{%
    \FormulaRefAuto{a=b\vdash b=a}[thm]{%
    \FormulaRefAuto{PowerSemigroupProductDef}[def]{1,2,4}}}
\end{tabproofwide}

\clearpage
\PowerInnerProofStatement{PowerSemigroupAssociativity}
\begingroup
\setlength{\LTpre}{0pt}
\begin{tabproofwide}
  \proofstepwidestar[1]{\SemigroupAlg{A}{\star}}{\rA}
  \proofstepwidestar[2]{%
    \hypertarget{power.proof.exists-congruence}{}%
    X,Y,Z\in\PowNE{A}}{\rA}
  \proofstepwidestar[3]{x\in X\land(y\in Y\land z\in Z)}{\rA}
  \proofstepwidestar[1,2,3]{%
    (x\star y)\star z=x\star(y\star z)}{%
    \FormulaRefAuto{SemigroupAssociativityAxiom}[ax]{1,%
    \rAIRepeat{3}{\FormulaRefAuto{PowerSemigroupCarrierElement}[thm]}{2,3}}}
  \proofstepwide[1,2]{(X\star_{\mathcal P}Y)\star_{\mathcal P}Z}{=}{%
    \begin{aligned}[t]
      &\{u\star z\mid u\in X\star_{\mathcal P}Y,\\[-2pt]
      &\qquad z\in Z\}
    \end{aligned}}{%
    \FormulaRefAuto{PowerSemigroupLeftOuterProductRepresentation}[thm]{1,2,2,2}}
  \proofstepwide[1,2]{}{=}{%
    \begin{aligned}[t]
      &\{(x\star y)\star z\mid\\[-2pt]
      &\qquad x\in X,y\in Y,z\in Z\}
    \end{aligned}}{%
    \FormulaRefAuto{TermImageCompositionLeftFromEquality}[thm]{%
    \FormulaRefAuto{PowerSemigroupProductDef}[def]{1,2,2}}}
  \proofstepwide[1,2]{}{=}{%
    \begin{aligned}[t]
      &\{x\star(y\star z)\mid\\[-2pt]
      &\qquad x\in X,y\in Y,z\in Z\}
    \end{aligned}}{%
    \FormulaRefAuto{TripleTermImageEqualityIntroductionRule}[thm]{%
    \mathord{[3]\vdots 4}}}
  \proofstepwide[1,2]{}{=}{%
    \begin{aligned}[t]
      &\{x\star v\mid x\in X,\\[-2pt]
      &\qquad v\in Y\star_{\mathcal P}Z\}
    \end{aligned}}{%
    \FormulaRefAuto{TermImageCompositionRightFromEquality}[thm]{%
    \FormulaRefAuto{PowerSemigroupProductDef}[def]{1,2,2}}}
  \proofstepwide[1,2]{}{=}{X\star_{\mathcal P}(Y\star_{\mathcal P}Z)}{%
    \FormulaRefAuto{PowerSemigroupRightOuterProductRepresentation}[thm]{1,2,2,2}}
  \proofstepwidestar[1,2]{\begin{aligned}[t]
    &(X\star_{\mathcal P}Y)\star_{\mathcal P}Z\\[-2pt]
    &\qquad=X\star_{\mathcal P}(Y\star_{\mathcal P}Z)
  \end{aligned}}{\rChain{5,9}}
\end{tabproofwide}
\endgroup

\par\Needspace{25\baselineskip}
\PowerProofStatement{NonemptyPowerSemigroup}
\begin{tabproofwide}
  \proofstepwidestar[1]{\SemigroupAlg{A}{\star}}{\rA}
  \proofstepwidestar[2]{X,Y\in\PowNE{A}}{\rA}
  \proofstepwidestar[3]{Z\in\PowNE{A}}{\rA}
  \proofstepwidestar[1,2]{X\star_{\mathcal P}Y\in\PowNE{A}}{%
    \FormulaRefAuto{PowerSemigroupProductClosure}[thm]{1,2,2}}
  \proofstepwidestar[1,2,3]{\begin{aligned}[t]
    &(X\star_{\mathcal P}Y)\star_{\mathcal P}Z\\[-2pt]
    &\qquad=X\star_{\mathcal P}(Y\star_{\mathcal P}Z)
  \end{aligned}}{\FormulaRefAuto{PowerSemigroupAssociativity}[thm]{1,2,2,3}}
  \proofstepwidestar[1]{%
    \hypertarget{power.proof.conclusion}{}%
    \SemigroupAlg{\PowNE{A}}{\star_{\mathcal P}}}{\begin{gathered}
    \FormulaRefAuto{SemigroupDef}[def](\\[-2pt]
    [2]\vdots4,\\[-2pt]
    [2,3]\vdots5)
  \end{gathered}}
\end{tabproofwide}

\par\Needspace{22\baselineskip}
\PowerProofStatement{PowerSemigroupProductMembership}
\begin{tabproofwide}
  \proofstepwidestar[1]{\SemigroupAlg{A}{\star}}{\rA}
  \proofstepwidestar[2]{X,Y\in\PowNE{A}}{\rA}
  \proofstepwidestar[1,2]{%
    \hypertarget{power.proof.PowerSemigroupProductWitnesses}{}%
    \hypertarget{power.proof.product-witness-form}{}%
    \begin{aligned}[t]
      &z\in X\star_{\mathcal P}Y\leftrightarrow{}\\[-2pt]
      &\quad\exists x\in X\,\exists y\in Y\;(z=x\star y)
    \end{aligned}}{%
    \rIEBack{\FormulaRefAuto{PowerSemigroupProductDef}[def]{1,2,2},%
    \FormulaRefAuto{B27TermImageMembership}[thm]}}
  \proofstepwidestar[1,2]{z\in X\star_{\mathcal P}Y\rightarrow z\in A}{%
    \rUE{\FormulaRefAuto{A\subseteq B\coloneq\forall x\,(x\in A\rightarrow x\in B)}[def]{%
    \FormulaRefAuto{PowerSemigroupProductSubset}[thm]{1,2,2}}}}
  \proofstepwidestar[1,2]{\begin{aligned}[t]
    &z\in X\star_{\mathcal P}Y\leftrightarrow{}\\[-2pt]
    &\quad\bigl(z\in A\land z\in X\star_{\mathcal P}Y\bigr)
  \end{aligned}}{%
    \FormulaRefAuto{P\rightarrow Q\vdash P\leftrightarrow(Q\land P)}[thm]{4}}
  \proofstepwidestar[1,2]{\begin{aligned}[t]
    &z\in X\star_{\mathcal P}Y\leftrightarrow{}\\[-2pt]
    &\quad\Bigl(z\in A\land{}\\[-2pt]
    &\qquad\exists x\in X\,\exists y\in Y\;(z=x\star y)\Bigr)
  \end{aligned}}{\rLRS{3,5}}
\end{tabproofwide}

In Zeile~5 wird der Satz aus Band~2 auf \(P\coloneqq z\in X\star_{\mathcal P}Y\)
und \(Q\coloneqq z\in A\) angewandt. Erst Zeile~6 ersetzt mit Zeile~3
das rechte Vorkommen von \(P\) durch die Existenzaussage.
